% Source: Wald GR, physical PDF pp. 483--485
%         (printed pp. 470--472; p. 472 is intentionally blank).
% Coverage: Appendix F, Units and Dimensions; Table F.1.
% Figures/tables/footnotes: no figures/footnotes; Table F.1.
% Status: source-reviewed and compile-clean.
% Uncertainties: none.

\section{Units and Dimensions}

\subsection*{Geometrized Units}

In this book, we have used \lq\lq geometrized units,\rq\rq\ where the gravitational
constant \(G\), and speed of light \(c\), are set equal to one. All quantities which
in ordinary units have dimension expressible in terms of length \(L\), time \(T\),
and mass \(M\), are given the dimension of a power of length in geometrized units.
Since \(G=c=1\), all factors involving \(G\) and \(c\) may be omitted from formulas,
and, indeed this is why it is convenient to use geometrized units in general
relativity. However, if one wishes to evaluate quantities in ordinary,
\lq\lq nongeometrized\rq\rq\ units, the factors of \(G\) and \(c\) must be reinserted.
This is easily done as follows.

In \lq\lq nongeometrized\rq\rq\ units, the dimension of \(c\) is \(L/T\) and the
dimension of \(G/c^2\) is \(L/M\). Hence, the \lq\lq conversion factor\rq\rq\
relative to geometrized units for a quantity with dimension of time is \(c\), while
the conversion factor for a quantity with dimension of mass is \(G/c^2\). More
generally, a quantity with dimension \(L^nT^mM^p\) in ordinary units has dimension
\(L^{n+m+p}\) in geometrized units and the conversion factor is \(c^m(G/c^2)^p\).
The dimensions and conversion factors for some frequently encountered quantities
are given in Table~\ref{tab:F.1}.

In order to convert a formula written in geometrized units to one which is valid in
nongeometrized units, one first must identify the nongeometrized dimension of all
quantities appearing in the equation. Then one simply obtains the conversion factor
for each quantity from Table~\ref{tab:F.1} or computes it by the above formula. If
one then \emph{multiplies} each quantity appearing in the equation by its conversion
factor, the resulting equation will be valid in nongeometrized units.

\subsection*{Planck Units}

For calculations involving quantum effects in general relativity, it is natural to
employ Planck units where \(\hbar\) is set equal to 1 in addition to \(G=c=1\). In
Planck units, all quantities which in ordinary units have dimension expressible in
terms of \(L\), \(T\), and \(M\) now become dimensionless. In particular, all lengths
are expressed as dimensionless multiples of the Planck length,
\(\ell_p=(G\hbar/c^3)^{1/2}\). To convert a formula valid in Planck units to one
valid in ordinary units, we simply identify the nongeometrized dimension of all
quantities appearing in the equation. Then we \emph{multiply} each such quantity by
its conversion factor, which equals its conversion factor for geometrized units
divided by \(\ell_p^n\), where \(L^n\) is the geometrized dimension of the quantity.

\begingroup
\renewcommand{\thetable}{F.1}
\begin{table}[H]
  \centering
  \caption{Conversion Factors to Geometrized Units}
  \label{tab:F.1}
  \small
  \renewcommand{\arraystretch}{1.2}
  \setlength{\tabcolsep}{9pt}
  \begin{tabular}{@{}lccc@{}}
    \toprule
    \textsc{Quantity} & \begin{tabular}{@{}c@{}}\textsc{Nongeometrized}\\
      \textsc{Dimension}\end{tabular} &
      \begin{tabular}{@{}c@{}}\textsc{Geometrized}\\\textsc{Dimension}\end{tabular} &
      \begin{tabular}{@{}c@{}}\textsc{Conversion}\\\textsc{Factor}\end{tabular}\\
    \midrule
    Acceleration & \(LT^{-2}\) & \(L^{-1}\) & \(c^{-2}\)\\
    Angular momentum & \(L^2T^{-1}M\) & \(L^2\) & \(G/c^3\)\\
    Electric charge (cgs) & \(L^{3/2}T^{-1}M^{1/2}\) & \(L\) & \(G^{1/2}/c^2\)\\
    Energy & \(L^2T^{-2}M\) & \(L\) & \(G/c^4\)\\
    Energy density & \(L^{-1}T^{-2}M\) & \(L^{-2}\) & \(G/c^4\)\\
    Force & \(LT^{-2}M\) & \(1\) & \(G/c^4\)\\
    Length & \(L\) & \(L\) & \(1\)\\
    Mass & \(M\) & \(L\) & \(G/c^2\)\\
    Mass density & \(L^{-3}M\) & \(L^{-2}\) & \(G/c^2\)\\
    Pressure & \(L^{-1}T^{-2}M\) & \(L^{-2}\) & \(G/c^4\)\\
    Time & \(T\) & \(L\) & \(c\)\\
    Velocity & \(LT^{-1}\) & \(1\) & \(c^{-1}\)\\
    \bottomrule
  \end{tabular}
\end{table}
\endgroup

In Planck units, the fundamental scales of length, time, mass, and other quantities
(with respect to which all physical quantities are expressed as dimensionless ratios)
are just the inverses of the above conversion factors. For the convenience of the
reader, we list some of these fundamental scales below, together with their values
in cgs units, calculated using the values \(c=3.00\times10^{10}\,
\mathrm{cm\ s}^{-1}\), \(G=6.67\times10^{-8}\,\mathrm{cm^3\ g^{-1}\ s^{-2}}\), and
\(\hbar=1.05\times10^{-27}\,\mathrm{erg\!-\!sec}\).

\begin{align*}
  \text{length:}\quad &\ell_p=(G\hbar/c^3)^{1/2}\approx1.6\times10^{-33}\,\mathrm{cm},\\
  \text{time:}\quad &t_p=\ell_p/c\approx5.4\times10^{-44}\,\mathrm{s},\\
  \text{mass:}\quad &m_p=\ell_pc^2/G\approx2.2\times10^{-5}\,\mathrm{g},\\
  \text{energy:}\quad &E_p=\ell_pc^4/G\approx2.0\times10^{16}\,\mathrm{ergs}
    \approx1.3\times10^{19}\,\mathrm{GeV},\\
  \text{mass density:}\quad &\rho_p=\ell_p^{-2}c^2/G
    \approx5.2\times10^{93}\,\mathrm{g\ cm}^{-3},\\
  \text{temperature:}\quad &T_p=E_p/k=\ell_pc^4/Gk
    \approx1.4\times10^{32}\,\mathrm{K}.
\end{align*}

% Source physical p. 485 (printed p. 472) is intentionally blank before References.
\clearpage
\thispagestyle{empty}
\null
\clearpage
